\begin{center}\begin{tikzpicture}[font=\small,>=stealth,line join=round]
\draw[gray!40,rounded corners=3pt] (0,0) rectangle (14.2,8.5);\node at (7.1,7.9) {易错五：偶函数在零处的值，可以不是零};
\node at (3.5,6.85){奇函数 $y=x^3$};\node at (10.6,6.85){偶函数 $y=x^2+2$};\begin{scope}[shift={(3.5,3.1)},x=0.95cm,y=0.75cm]\draw[->,gray] (-2.2,0)--(2.2,0) node[right]{$x$};\draw[->,gray] (0,-1.5)--(0,4.1) node[above]{$y$};\draw[blue,thick,domain=-1.1:1.55,samples=150] plot(\x,{(\x)^3});\fill[orange!85!black] (0,0) circle(2.2pt) node[below right=5pt] {$(0,0)$};\end{scope}\begin{scope}[shift={(10.5,3.1)},x=0.95cm,y=0.75cm]\draw[->,gray] (-2.2,0)--(2.2,0) node[right]{$x$};\draw[->,gray] (0,-1.5)--(0,4.1) node[above]{$y$};\draw[blue,thick,domain=-1.4:1.4,samples=150] plot(\x,{(\x)^2+2});\fill[orange!85!black] (0,2) circle(2.2pt) node[below right=5pt] {$(0,2)$};\fill[orange!85!black] (-1,3) circle(2.2pt);\fill[orange!85!black] (1,3) circle(2.2pt);\draw[gray,dashed] (-1,3)--(1,3);\end{scope}
\node[text width=12.8cm,align=center] at (7.1,.65) {两图同尺度；左图要求相反数配对，右图只要求左右高度相同。};
\end{tikzpicture}\end{center}
